\section{An exact-rational certificate for a simple critical-line zero}
\label{sec:example}

The preceding general contour bounds establish totality but are far too
large for a practical run. The archive therefore also contains a small
independent calculation using much sharper local information. This is a
worked arithmetic certificate, with its analytic interpretation proved
below. It illustrates the distinction between executing a calculation
and formalizing that calculation's soundness in the object theory.

Take the exact rational data
\begin{equation}
 c=\frac12+14.1347251417347\,\ii,
 \qquad r=h=10^{-6},\qquad B_2=4096.
 \label{eq:example-data}
\end{equation}
The decimal in this display is a terminating rational, not an assumed
exact zeta ordinate. Use~\eqref{eq:em} with $N=64$ and $m=10$ to
enclose $f(c)$, $f(c-h)$ and $f(c+h)$, where $f=\zeta$. The
implementation uses exact rational arithmetic, 384-bit dyadic centres,
outward-rounded radii, the atanh logarithm series, and the exponential
Taylor series. The Bernoulli numbers are computed by their rational
recurrence. The remainder is bounded by
\eqref{eq:em-bound}, multiplied by two for the $\normone{\cdot}$ radius.

Let $d$ be the rational centre of
$(f(c+h)-f(c-h))/(2h)$, and let $e_d$ be its evaluation error plus
$B_2h/2$. Then
\begin{equation}
 |f'(c)-d|\le e_d.
 \label{eq:example-derivative}
\end{equation}
To see this, integrate $f'$ over the real segment $[c-h,c+h]$.
The mean error from replacing $f'(c+x)$ by $f'(c)$ is at most
$(2h)^{-1}\int_{-h}^h B_2|x|\,dx=B_2h/2$.

Here is a completely elementary justification of the second-derivative
bound. On the disk and the finite-difference segment,
$\sigma\ge1/4$, $|s|\le16$ and $|s-1|\ge1$.
Differentiating~\eqref{eq:euler-basic} twice and using
\eqref{eq:logmoments} gives
\[
 |\zeta''(s)|\le 2+2\cdot16+16\cdot128=2082<4096.
\]
The same inequalities hold throughout a neighborhood of these sets,
as can be checked directly from~\eqref{eq:example-data}.

The script verifies, as exact rational comparisons,
\begin{align}
 A_0+e_dr+\tfrac12B_2r^2&< |d|_\infty r,
       \label{eq:example-rouche}\\
 |d|_\infty-e_d-B_2r&>0,
       \label{eq:example-simple}
\end{align}
where $A_0$ is the certified upper bound for $|f(c)|$.
Some outward numerical summaries of the exact inequalities are
\begin{center}
\begin{tabular}{ll}
\toprule
Quantity & Conservative decimal summary\\
\midrule
$A_0$ & less than $5.639\times10^{-15}$\\
$\Re d$ & between $0.78329651186710$ and $0.78329651186712$\\
$\Im d$ & between $0.12469982974821$ and $0.12469982974823$\\
Left side of~\eqref{eq:example-rouche}
  & less than $4.097\times10^{-9}$\\
Right side of~\eqref{eq:example-rouche}
  & greater than $7.832\times10^{-7}$\\
Lower bound in~\eqref{eq:example-simple}
  & greater than $0.77715$\\
\bottomrule
\end{tabular}
\end{center}
The JSON file retains the exact rational values; these displayed
decimals are not used to make the decisions.

Taylor's formula with the explicit integral remainder gives, on
$|z-c|=r$,
\[
 |f(z)-d(z-c)|\le A_0+e_dr+B_2r^2/2<|d(z-c)|.
\]
Rouch\'e stability therefore gives count one in the disk. Also
$|f'(z)|\ge |d|_\infty-e_d-B_2r>0$ throughout the disk, so its
zero is simple.

\subsection{A primitive recursive point name}
There is an even more explicit constructive existence argument here.
Use the fixed rational $d\ne0$ and the map
\[
             \Phi(z)=z-f(z)/d.
\]
The rational certificate checks
\begin{equation}
 \kappa:=\frac{e_d+B_2r}{|d|_\infty}<\frac1{64},\qquad
 \frac{A_0}{|d|_\infty}<\frac r{64}.
 \label{eq:example-contraction}
\end{equation}
Thus $\Phi$ maps the closed disk into the concentric disk of radius
$r/32$ and is a contraction of factor at most $1/64$.
Starting at $c$, $n$ exact iterations have successive errors bounded
geometrically, so their distance from the root is at most
$2r64^{-n}$. For a requested precision, perform a predetermined
number of iterations and evaluate each operation with a predetermined
rounding budget; the geometric stability bound controls the sum of
rounding errors. For example, use $n=p+30$ iterations and a per-step
error at most $2^{-p-40}r$, increasing the elementary working precision
by the explicit division bound $|d|>1/2$. The final error is less
than $2^{-p}$. This is a primitive recursive Cauchy name. The same
contraction estimate proves uniqueness.

The disk is invariant under $z\mapsto1-\bar z$. The completed
functional equation sends its root to another root in the same disk;
uniqueness implies $\rho=1-\bar\rho$, hence $\Re\rho=1/2$.
The certificate therefore proves one simple critical-line zero in
the disk~\eqref{eq:example-data}. It does not claim that this is the
lowest positive zero; that extra statement would require a count in
the entire lower region.

\subsection{What was executed}
\code{python3 code/rational_zeta.py} recomputes all three zeta
enclosures and all displayed strict rational tests. It writes
\code{verification/simple_zero.json}. No floating-point special-function
evaluation, zero table or root finder is used by this program.
Its correctness is justified by the mathematical estimates just given.
The program is not a proof checker for full $\ipra$, and its successful
execution is not reported as a Lean kernel proof.
